\chapter{File boundaries and binary images}
\MARSChapterHeading{File boundaries and binary images}
\label{revision:files}

A framebuffer is useful while the program is running. A file lets another program
recover the image later, but the file must describe more than its pixel values.
The reader needs dimensions, a sample representation and an exact boundary
between the description and the pixels. It also needs to distinguish malformed
content from a failed attempt to read that content.

The complete \texttt{ppm-codec.asm} in \texttt{revision/labs/files} reads,
decodes, serializes and saves a small binary PPM image. PPM means Portable Pixel
Map. We use a documented subset of its P6 form rather than implementing every
permitted PPM variant. This chapter assumes byte addressing, bounded parsing,
RGB words and nested function calls. Its independent task produces a complementary color image while preserving the source framebuffer.

\section{Know where text ends and samples begin}

The default two-by-two file begins with eleven ASCII header bytes:
\begin{verbatim}
P6\n2 2\n255\n
\end{verbatim}
Here \texttt{\textbackslash n} denotes one LF byte, hexadecimal \texttt{0a};
it is not a backslash followed by the letter n in the file. \texttt{P6} identifies
the binary form. Decimal width and height are followed by Maxval 255, meaning
that each component uses one byte with values zero through 255. The remaining
twelve bytes describe four pixels, each as red, green and blue, row by row from
top left to bottom right.

\begin{figure}[htbp]
\centering
\begin{tabular}{r l l}
\toprule
Offset & Bytes & Interpretation\\
\midrule
0--2 & \texttt{50 36 0a} & \texttt{P6} and LF\\
3--6 & \texttt{32 20 32 0a} & width two, space, height two, LF\\
7--10 & \texttt{32 35 35 0a} & Maxval 255 and final LF\\
11--13 & \texttt{ff 00 00} & top-left red pixel\\
14--16 & \texttt{00 ff 00} & top-right green pixel\\
17--19 & \texttt{00 00 ff} & bottom-left blue pixel\\
20--22 & \texttt{ff ff ff} & bottom-right white pixel\\
\bottomrule
\end{tabular}
\caption{The exact 23-byte worked P6 image. Header bytes have ASCII meaning; at offset eleven interpretation changes to RGB samples. No row delimiter or terminating zero follows the raster.}
\label{fig:revision-ppm-layout}
\end{figure}


A header gap may contain ASCII space, bytes nine through thirteen, or a
\texttt{\#} comment terminated by LF. Gaps are required after the magic and
between numeric tokens. Our subset requires positive decimal dimensions from
one through 64 without leading zeros, Maxval exactly 255, and exactly one LF
immediately after that token. It accepts one image, a header no longer than
256 bytes and no bytes after the raster. Other valid PPM variants may fall
outside these local rules.

After that final LF, all bytes are raster samples. A newline, space, zero or
\texttt{\#} has no lexical role there. A one-by-one file can end in bytes
\texttt{0a 23 00}, representing red ten, green 35 and blue zero. A generic
whitespace/comment skipper called after Maxval would eat part of that pixel.
The parser must change interpretation at the specified boundary.

The parser returns a sixteen-byte descriptor: unsigned word width at offset
zero, unsigned word height at four, and an unsigned doubleword raster offset
at eight. The offset is measured from the supplied file-buffer base. It is not
an absolute pointer. Another caller may supply the same bytes at a different
address without changing the descriptor's offset.

\section{Prove the complete extent before publishing metadata}

\texttt{ppm\_parse} receives full byte pointer \texttt{\$a0}, unsigned selected
byte count \texttt{\$a1}, and an eight-byte-aligned metadata pointer
\texttt{\$a2}. It rejects a count above 12544 before pointer access. This bound
is 256 possible header bytes plus $64\cdot64\cdot3$ raster bytes. Empty input
is malformed, not a successful zero-pixel image.

The scan limit is the smaller of selected count and 256. Every header byte load
must be below that limit, including bytes inside a comment or decimal token.
The decimal accumulator starts only on a digit one through nine; subsequent
digits use $\mathit{value}=10\cdot\mathit{value}+\mathit{digit}$. The worked parser rejects an accumulator
above 255 immediately, then applies the smaller dimension bound where required.
This avoids accepting a long overflowing decimal string through wraparound.

Let $H$ be the position immediately after the required final LF. The parser has
proved $H\leq \mathit{selectedCount}$ and $H\leq256$. For accepted width $W$ and height
$T$, require
\[
  \mathit{selectedCount}-H=3WT.
\]
The product is bounded by 12288, and the subtraction is safe because the scan
already established its ordering. Equality rejects both a missing sample and
extra trailing data. Mapped padding does not make either selected extent valid.

Only after this check does the parser publish the descriptor:
\Needspace{10\baselineskip}
\MARSCodeLabel{Publish metadata after the complete extent check}
\begin{minted}{mars-idea}
sw $t6, 0($a2)
sw $t7, 4($a2)
sd $t0, 8($a2)
move $v0, $zero
jrc $ra
nop
\end{minted}
Here \texttt{\$t6/\$t7} are accepted width/height, \texttt{\$t0} is $H$, and
\texttt{\$v1} already holds the pixel count. A failure returns status one for
malformed data or two for count bound, with \texttt{\$v1}=0 and metadata
unchanged. Writing plausible dimensions early would violate that guarantee if
a later raster-length check failed.

\section{Convert sample order rather than copying words}

\texttt{ppm\_decode} validates into a private descriptor before checking the
unsigned pixel capacity in \texttt{\$a3}. Malformed content therefore wins over
insufficient capacity. On success it publishes one framebuffer word
\texttt{0x00RRGGBB} per pixel. The framebuffer pointer is four-byte aligned;
file bytes may begin at an arbitrary byte address.

File triplet \texttt{12 34 56} becomes word \texttt{00123456}. In our
little-endian memory that word occupies bytes \texttt{56 34 12 00}.
A four-byte copy from the file would mix neighboring pixel data and confuse
file component order with memory word order. The worked loop loads exactly
three unsigned bytes, shifts red by sixteen and green by eight, combines the
components and stores one word.

\begin{figure}[htbp]
  \centering
  \begin{tabular}{l l}
    \toprule
    Representation                  & Example              \\
    \midrule
    File RGB samples                & \texttt{12 34 56}    \\
    Framebuffer word                & \texttt{00123456}    \\
    Little-endian framebuffer bytes & \texttt{56 34 12 00} \\
    Word with ignored high byte     & \texttt{aa123456}    \\
    Complementary file samples      & \texttt{ed cb a9}    \\
    \bottomrule
  \end{tabular}
  \caption{Component order in the file is independent of byte order within a machine word. Encoding ignores the high byte; the independent negative-image encoder complements the three color components without changing that source word. All displayed values are hexadecimal.}
  \label{fig:revision-ppm-order}
\end{figure}


The default raster is red, green, blue and white. Decoded words are
\texttt{00ff0000}, \texttt{0000ff00}, \texttt{000000ff}, \texttt{00ffffff}.
The decoder returns dimensions packed in \texttt{\$v1}: height in the high
32 bits and width in the low 32. This differs from the parser's pixel count.
For two by two, the packed value is hexadecimal \texttt{0000000200000002},
not four. Keep the units and representation of each result beside its interface.

The encoder takes this packed dimension value in \texttt{\$a1}, a framebuffer
pointer in \texttt{\$a0}, byte destination in \texttt{\$a2}, and unsigned byte
capacity in \texttt{\$a3}. It validates both dimensions before capacity, and
both before accessing caller memory. It emits the canonical header shown above,
followed by RGB triplets, ignoring each input word's high byte. No terminator
belongs to the file.

For these bounded dimensions, header length is eleven plus one for each
two-digit dimension. Total output length is therefore
$3WT+11+[W\geq10]+[T\geq10]$, where a bracket contributes one when its
condition holds. Width nine and height ten need a twelve-byte header. The
length proof permits a complete capacity check before the first store. The
encoder returns exact byte length in \texttt{\$v1}; all its failures return
zero length and leave output unchanged.

\section{Use a descriptor to continue a file operation}

MARS supplies educational file services through \texttt{syscall}. Put a service
number in \texttt{\$v0}, prepare its arguments, then execute the instruction.
These service numbers belong to the simulator's environment; they are not a
portable hardware or operating-system interface.

\begin{center}
  \begin{tabular}{r l l}
    \toprule
    Service  & Inputs                                   & Result                                           \\
    \midrule
    13 Open  & \texttt{\$a0} path, \texttt{\$a1} flags  & descriptor or negative failure                   \\
    14 Read  & descriptor, destination, requested bytes & transferred count, zero EOF, or negative failure \\
    15 Write & descriptor, source, requested bytes      & accepted count or negative failure               \\
    16 Close & \texttt{\$a0} descriptor                 & no count result used here                        \\
    \bottomrule
  \end{tabular}
\end{center}
For read/write, arguments occupy \texttt{\$a0/\$a1/\$a2} in the order shown.
Paths are NUL-terminated strings. Their termination rule does not apply to raster
bytes or serialized file contents. Open flag zero selects reading; flag one
selects writing with truncation. Relative paths are resolved from the configured
working directory, which the lab makes explicit.

A successful open returns a small integer descriptor. Retain that returned
number; do not assume it is always three or confuse it with the path pointer.
The program owns the descriptor until it closes it. Every path after a successful
open, including read/write failure, must close exactly that descriptor once.
A failed open returns no descriptor to close.

The read/write result overwrites \texttt{\$v0}, so reload the service number
before the next call. Keep live pointers, totals and descriptor in preserved
registers or private storage across helper calls. A returned transfer count is
progress within the current file operation, not the file's total length.

\section{Read through EOF, including at an exact capacity boundary}

\texttt{ppm\_load} receives path, byte buffer, unsigned byte capacity and metadata
pointer in \texttt{\$a0} through \texttt{\$a3}. It accepts capacity one through
12544 before accessing paths, buffers or files. Its byte buffer may contain a
prefix on failure, but metadata remains unchanged until complete validation.

Keep $used$ bytes already read. Request $capacity-used$ bytes into
$buffer+used$. A positive result advances $used$ by that result, even if smaller
than the request. A short read is not EOF: continue until a read returns zero.
A negative result is I/O failure. The worked loop's request preparation is:
\Needspace{11\baselineskip}
\MARSCodeLabel{Read into the unfilled suffix}
\begin{minted}{mars-idea}
move $a0, $s3
daddu $a1, $s0, $s4
dsubu $a2, $s1, $s4
li $v0, 14
syscall
\end{minted}
The saved state holds descriptor in \texttt{\$s3}, buffer base in
\texttt{\$s0}, capacity in \texttt{\$s1}, and used count in \texttt{\$s4}.
The surrounding loop checks errors, zero and the updated capacity invariant.

When $used=capacity$, another buffer read would overrun the selected destination,
but reaching capacity alone cannot establish EOF. The function requests one
extra byte into its private stack probe. Zero means the file ended exactly at
capacity. A positive result means the file is too large and returns bound status
two; a negative result is I/O status three. The probe byte never goes into caller
padding. This distinguishes an exactly full valid file from its one-byte-longer
counterexample.

After EOF, close the descriptor, then call the parser on exactly $used$ bytes.
Only successful parsing publishes metadata. Success returns file byte count in
\texttt{\$v1}. A new open starts again at the beginning; no hidden descriptor or
cursor survives a call. The lab tests these rules with short transfers and errors
as well as ordinary host files.

\section{Validate before truncation, then finish short writes}

\texttt{ppm\_save} takes path, serialized byte pointer and unsigned selected
count in \texttt{\$a0/\$a1/\$a2}. It checks the bound and complete format
before opening the output in truncate mode. Otherwise malformed data could destroy
an existing good file before the function discovered its error.

Once open succeeds, keep a count of bytes already accepted. Request the remaining
suffix, advance only by the positive returned count, and repeat until all selected
bytes have been accepted. Treat a zero write as failure to make progress, rather
than looping forever. Close the descriptor on both completion and failure.
Success returns the count; I/O failure returns status three and zero count.

This is not a transaction on the host filesystem. A failure after truncation may
leave a partial file. The stronger guarantee is that format rejection occurs
before opening, and that the function neither invents a successful full length
nor leaks its descriptor after an I/O error. A higher-level application that
needs replacement as one operation would need a different file protocol.

The current MARS write service reads one byte beyond its requested memory range
while preparing its host buffer. Our serialized allocations retain a readable
guard after the selected bytes for that implementation behavior. The guard is
not part of PPM, not a terminator to serialize, and not included in the requested
or returned file length. The selected file count remains authoritative. This
requirement is specific to the current service implementation.

\section{Join verified functions into a complete artifact}

The worked caller's mode zero parses its embedded image without file access.
Modes one and two decode and encode; three and four load and save. Mode five
composes load, decode, encode and save. It passes the loader's byte count to the
decoder, the decoder's packed dimensions to the encoder, and the encoder's byte
count to the saver. An error stops that chain before a later stage can interpret
an error result as a size or dimension value.

The public \texttt{tools/run\_revision\_ppm.py} makes a fresh working directory,
creates a deterministic input, selects pipeline mode in a separate source copy,
runs the actual MARS CLI and verifies the host output's canonical bytes. It also
starts with a longer old output to check truncation. A commented input may
produce a shorter canonical header while retaining the same samples. A byte-for-
byte comparison with the original file would then test the wrong property.

\begin{MARSWorkedExample}{Four colors through the complete pipeline}
  The input has eleven header bytes and twelve raster bytes. Load returns 23,
  parse describes width two, height two and offset eleven, decode produces four
  words and packed dimensions, encode returns 23, and save accepts those 23 bytes.
  The output is a complete image file. None of the path terminators, framebuffer
  high bytes, backing padding or write guard belongs to that file.
\end{MARSWorkedExample}

All functions preserve full saved registers, frame pointer, stack pointer and
return path, and perform no console I/O. Input and output objects are disjoint.
Status zero means success, one malformed content or invalid dimensions, two a
count/capacity bound, and three a file I/O failure. Failure length is zero, with
the publication limits already stated for each interface. The lab's public
checks test those limits independently of the worked register allocation.

\section{Produce complementary colors independently}

Complete \texttt{ppm\_encode\_negative} in \texttt{ppm-negative.asm}. Its arguments,
dimension range, canonical header, count/capacity rules and failure publication
match \texttt{ppm\_encode}. For each source channel $x$, serialize $255-x$;
ignore the high byte and leave every source word unchanged. Red becomes cyan,
green becomes magenta, blue becomes yellow and white becomes black.

You may reuse the supplied encoder and other helpers. One design changes the
sample calculation while emitting the raster; another encodes successfully then
transforms only its raster bytes. Explain how each design keeps the header
unchanged, locates the exact raster, preserves source and saved state, and
performs no output writes on failure. The accepted output interval may be
written during a successful computation; an unchanged final buffer alone does
not establish a failure guarantee.

Use a high-byte value other than zero, dimensions across the nine/ten boundary,
exact and insufficient capacity, and a later successful call after failure.
The independent task includes a complete host-file run using your implementation,
not only an inspection of its in-memory prefix. The written lab supplies the
checks and commands without prescribing a frame or instruction schedule.

\section*{Sources and laboratory connection}
The \href{https://netpbm.sourceforge.net/doc/ppm.html}{Netpbm PPM specification}
defines the P6 header and raster model. The local rules in
\path{packages/029-binary-ppm/} deliberately restrict dimensions, numeric spelling,
Maxval, final separator and image count. The educational file ABI and guard
behavior are checked against the pinned MARS syscall implementations. The lab
connects source-byte boundaries, descriptor publication, actual host artifacts,
short-transfer reasoning and the independent complementary-color encoder.
